\documentclass[10pt,a4paper]{amsart}
\usepackage[margin=19mm]{geometry}
\usepackage{amsmath,amssymb,mathtools}
\usepackage[scr=rsfs,cal=euler]{mathalfa}
\usepackage{xfrac}
\usepackage[unicode,hidelinks,bookmarks=false]{hyperref}
\newcommand{\ie}{i.e.\ }
\newcommand{\cf}{cf.\ }
\newcommand{\resp}{respectively\ }
\newcommand{\Subsection}{\subsection}
\providecommand{\nobreakdash}{\nobreak\hbox{-}}
\setlength{\emergencystretch}{2em}
\multlinegap0pt
\usepackage{amssymb}
\usepackage{xcolor}
\newtheorem{lemma}{Lemma}
\newtheorem{theorem}{Theorem}
\title{An improved volume bound under Ricci and scalar curvature lower bounds}
\author{Kwok-Kun Kwong}
\date{Source: arXiv:2608.19196v2 (CC BY 4.0)}
\begin{document}
\maketitle

\noindent\emph{Source and licence.} An improved volume bound under Ricci and scalar curvature lower bounds. Version arXiv:2608.19196v2 (CC BY 4.0), distributed by arXiv under the Creative Commons Attribution 4.0 International licence, http://creativecommons.org/licenses/by/4.0/, as recorded in the arXiv licence metadata for this exact version and shown on the abstract page https://arxiv.org/abs/2608.19196v2. Full licence notice: \texttt{licenses/arxiv-2608.19196-CC-BY-4.0.txt}.\par\smallskip

\noindent\textit{Editorial scope.} Counted unit: the article's theorem thm:continuum-shuffle (source0.tex lines 727--736). The article's complete original proof of it is delivered with it (source0.tex lines 738--803, one \texttt{proof} environment of 66 lines). No mathematics is written, repaired or re-derived: the statement and the proof are the article's own lines, in the article's own notation, and no retained line is substituted or re-flowed.  Retained uncounted as the source-local context the proof needs: lines 482--484; lines 488--494; lines 558--568; lines 635--654; lines 640--644.  Nothing was omitted from any retained span, so the package carries no omission record.  No retained line contains a citation, so the wrapper carries no bibliography and no bibliography entry.  No retained line contains a figure environment, an image-inclusion command or drawing code, so no image is delivered and the package declares no asset dependency.  Editorial formatting only, in addition: an AMS \texttt{amsart} preamble that restates the article's own declarations and macros used by the retained lines, namely the environments \texttt{lemma}, \texttt{theorem} on separate counters, together with the standard packages those lines need.  The retained environments print in this document as Lemma 1, Theorem 1 in order of appearance, whereas the article prints the same items with the numbering of its own full document.  The complete native source of the article as distributed in the arXiv e-print for this exact version is bundled unchanged as \texttt{sources/208069/source0.tex}.
\begin{equation}\label{eq:ordinary-scalar-jacobi}
u_q''+q(t)u_q=0, \qquad u_q(0)=0, \qquad u_q'(0)=1.
\end{equation}

\begin{equation}\label{eq:definition-stopped-solution}
j_q(t)=
\begin{cases}
u_q(t), & 0\le t\le \tau_q, \\
0, & t>\tau_q.
\end{cases}
\end{equation}

\begin{lemma}
\label{lem:stopped-continuity}
Let $q_\nu, q\in L^\infty[0, T]$ and suppose that
$$
\sup_\nu\lVert q_\nu\rVert_{L^\infty[0, T]}<\infty, \qquad \lVert q_\nu-q\rVert_{L^1[0, T]}\longrightarrow0.
$$
Let $j_{q_\nu}$ and $j_q$ be the stopped solutions defined by \eqref{eq:ordinary-scalar-jacobi} and \eqref{eq:definition-stopped-solution}. Then
\begin{equation*}\label{eq:stopped-continuity}
j_{q_\nu}(T)\longrightarrow j_q(T).
\end{equation*}
\end{lemma}

\begin{lemma}
\label{lem:finite-valued-integral}
Let $(X, m)$ be a probability measure space, and let $0=t_0<t_1<\cdots<t_{\mathrm{N}}=T$. Let $a_1, \ldots, a_s$ are distinct real numbers. Suppose we have a family of measurable functions
$$q_r:X\longrightarrow\{a_1, \ldots, a_s\}$$
for $r=1, \cdots, {\mathrm{N}}$, such that
\begin{equation}\label{eq:equal-level-measures}
m\{z:q_1(z)=a_i\}=
m\{z:q_2(z)=a_i\}=\cdots=
m\{z:q_{\mathrm{N}}(z)=a_i\} \text{ for each } a_i.
\end{equation}
For each $z\in X$, define
$$
\mathrm{p}_z(t)=q_r(z) \qquad\text{when }t_{r-1}<t<t_r,
$$
and let $j_z:=j_{\mathrm{p}_z}$ be its stopped solution. Then, for every $\mu\in \mathbb { N }$,
\begin{equation}\label{eq:finite-valued-integral}
\int_Xj_z(T)^\mu\, dm(z)
\le \int_X\widehat s_{q_1(z)}(T)^\mu\, dm(z).
\end{equation}
\end{lemma}

\begin{equation}
m\{z:q_1(z)=a_i\}=
m\{z:q_2(z)=a_i\}=\cdots=
m\{z:q_{\mathrm{N}}(z)=a_i\} \text{ for each } a_i.
\end{equation}

\begin{theorem} \label{thm:continuum-shuffle}
Let $X$ be a compact metric space equipped with a Borel probability measure $m$. Let $\Phi:[0, \infty)\times X\longrightarrow X$ be a continuous flow, and write $\Phi_t=\Phi(t, \cdot)$. Supppose $m(\Phi_t^{-1}(E))=m(E)$ for every Borel set $E\subset X$ and every $t\ge 0$. Let $K$ be a continuous function on $X$. For each $z\in X$, let $j_z$ be the stopped solution of
\begin{equation*}\label{eq:flow-jacobi}
j''(t)+K(\Phi_tz)j(t)=0, \qquad j(0)=0, \qquad j'(0)=1.
\end{equation*}
Then for every positive integer $\mu$ and every $t\ge 0$,
\begin{equation*}\label{eq:continuum-shuffle}
\int_Xj_z(t)^\mu\, dm(z) \le \int_X\widehat s_{K(z)}(t)^\mu\, dm(z).
\end{equation*}
\end{theorem}

\begin{proof}
Fix $T>0$. Since $K(X)$ is a compact subset of $\mathbb R$, there is a sequence of Borel functions $K_{\mathrm{N}}\colon X\to\mathbb R$, each taking values in a finite set, such that
\begin{equation}\label{eq:uniform-finite-valued-approximation}
\|K_{\mathrm{N}}-K\|_{L^\infty(X)}\longrightarrow0.
\end{equation}
In particular,
\begin{equation}\label{eq:uniform-coefficient-bound}
\sup_{\mathrm{N}}\|K_{\mathrm{N}}\|_{L^\infty(X)}+\|K\|_{L^\infty(X)}\le M
\end{equation}
for some $M<\infty$.

Let $t_r=\frac{rT}{{\mathrm{N}}}$, \;$0\le r\le {\mathrm{N}}$,
and define the piecewise constant coefficient
\begin{equation}\label{eq:discrete-flow-coefficient}
q_{{\mathrm{N}}, z}(t)=K_{\mathrm{N}}(\Phi_{t_{r-1}}z)
\qquad
\text{for }t_{r-1}<t<t_r.
\end{equation}
Let $j_{{\mathrm{N}}, z}$ be its stopped solution.

For every value $a$ taken by $K_{\mathrm{N}}$, measure preservation gives
\begin{align*}
m\{z:K_{\mathrm{N}}(\Phi_{t_{r-1}}z)=a\}
& =
m\left(\Phi_{t_{r-1}}^{-1}\{z:K_{\mathrm{N}}(z)=a\}\right)\\
& =
m\{z:K_{\mathrm{N}}(z)=a\}.
\end{align*}

Thus the coefficient functions $z\longmapsto K_{\mathrm{N}}(\Phi_{t_{r-1}}z)$ satisfy \eqref{eq:equal-level-measures}. Since $\Phi_0$ is the identity, Lemma \ref{lem:finite-valued-integral} gives
\begin{equation}\label{eq:approximating-integral-comparison}
\int_Xj_{{\mathrm{N}}, z}(T)^\mu\, dm(z)
\le \int_X\widehat s_{K_{\mathrm{N}}(z)}(T)^\mu\, dm(z).
\end{equation}

The map $(t, z)\longmapsto K(\Phi_tz)$ is uniformly continuous on $[0, T]\times X$. It follows
from \eqref{eq:uniform-finite-valued-approximation} and
\eqref{eq:discrete-flow-coefficient} that as $\mathrm{N}\to\infty$,
\begin{align*}
& \sup_{z\in X}\sup_{0\le t\le T} \left|q_{{\mathrm{N}}, z}(t)-K(\Phi_tz)\right|\\
\le& \lVert K_{\mathrm{N}}-K\rVert_{L^\infty(X)} + \sup_{\substack{z\in X, \;s, t\in[0, T] |s-t|\le T/{\mathrm{N}}}} \left|K(\Phi_sz)-K(\Phi_tz)\right| \longrightarrow0.
\end{align*}
Lemma \ref{lem:stopped-continuity} therefore implies that, for every $z\in X$, $\displaystyle \lim_{\mathrm{N}\to \infty}j_{{\mathrm{N}}, z}(T)= j_z(T)$.

Let $u_{{\mathrm{N}}, z}$ denote the ordinary, unstopped solution associated with $q_{{\mathrm{N}}, z}$. By \eqref{eq:uniform-coefficient-bound}, the corresponding first-order system has uniformly bounded coefficients. As in Lemma \ref{lem:stopped-continuity}, Gronwall's inequality gives a constant $C_T$, independent of ${\mathrm{N}}$ and $z$,
such that
$$
|u_{{\mathrm{N}}, z}(t)|+|u_{{\mathrm{N}}, z}'(t)|\le C_T
\qquad
(0\le t\le T).
$$
Since $j_{{\mathrm{N}}, z}$ agrees with $u_{{\mathrm{N}}, z}$ before its first positive zero and is zero thereafter,
$$
0\le j_{{\mathrm{N}}, z}(T)^\mu\le C_T^\mu.
$$
The same bound holds, after increasing $C_T$ if necessary, for $\widehat s_{K_{\mathrm{N}}(z)}(T)^\mu$. Moreover, as $\mathrm{N} \rightarrow \infty$,
$$
\widehat s_{K_{\mathrm{N}}(z)}(T) \longrightarrow \widehat s_{K(z)}(T),
$$
because $a\mapsto\widehat s_a(T)$ is continuous on $\mathbb R$. Applying dominated convergence theorem in \eqref{eq:approximating-integral-comparison} then gives
$$
\int_Xj_z(T)^\mu\, dm(z) \le \int_X\widehat s_{K(z)}(T)^\mu\, dm(z).
$$

Since $T>0$ was arbitrary, the theorem follows.
\end{proof}

\end{document}
